\section{Introduction}
\label{sec:introduction}

Urban traffic congestion represents an escalating challenge in metropolitan transportation management, directly correlating with economic inefficiencies, elevated greenhouse gas emissions, and life-threatening delays for emergency response vehicles (EVs)~\cite{webster1958, mnih2015human}. Conventional municipal intersections continue to rely predominantly on fixed-time or rule-based actuated signal controllers. While fixed-cycle timers are straightforward to implement, their inability to adapt to asymmetric flow rates leads to severe green starvation, where empty corridors are granted green allocations while congested lanes endure prolonged halts~\cite{lowrie1990scats}.

The limitations of traditional infrastructure become critically apparent during emergency vehicle routing. Ambulances, fire engines, and law enforcement units often face gridlocked intersection approaches. Existing preemption methods generally rely on either visual detection or acoustic siren triggers that force immediate all-red or green overrides~\cite{cucchiara2003detecting, salamon2014urbansound}. However, naively forcing preemption without modeling the surrounding traffic dynamics frequently causes severe secondary congestion shockwaves in opposing corridors. Once the emergency vehicle clears the intersection, cross-street queues reach saturation, and classical systems require excessive duration to recover equilibrium.

Deep Reinforcement Learning (DRL) has emerged as a promising paradigm for adaptive traffic signal control~\cite{wei2018intellilight, liang2019deep}. By formulating traffic actuation as a Markov Decision Process (MDP), reinforcement learning agents can autonomously learn non-linear policies mapping observed queue distributions to optimal phase allocations. Nevertheless, deploying pure DRL policies directly to safety-critical urban intersections introduces three severe operational challenges:
\begin{enumerate}
    \item \textbf{Absence of Physical Safety Bounds:} Pure end-to-end RL policies lack deterministic clearance bounds. An unconstrained agent can output erratic or extremely short phase durations (e.g., 2\,s green allocations during exploration), inducing high risks of rear-end vehicular collisions.
    \item \textbf{Post-Preemption Gridlock Cascades:} When an emergency override terminates, the RL state distribution experiences an abrupt out-of-distribution queue surge on cross-corridors, leading to destabilization or policy collapse.
    \item \textbf{Unrealistic Benchmark Baselines:} A prevalent shortcoming in contemporary literature is evaluating learning-based models solely against naive fixed-time controllers, masking whether the proposed architecture genuinely outperforms standard industrial heuristics such as Max-Pressure or gap-out actuated controllers.
\end{enumerate}

To overcome these critical limitations, this paper presents \textbf{AURA} (\textit{Adaptive Urban Routing and Actuation}), a decentralized reinforcement learning framework formulated as a \textbf{Semi-Markov Decision Process (SMDP)} with action-duration bounds for resilient traffic regulation and emergency corridor preemption. Instead of allowing the neural network to output arbitrary phase lengths, AURA delegates the discrete phase direction to a Deep Q-Network while mathematically bounding action durations: a statutory minimum green allocation of \textbf{10\,s} to ensure safe physical clearance, a queue-proportional dynamic scaling ceiling at \textbf{45\,s}, continuous-time queue dynamics modeled via \textbf{Neural Ordinary Differential Equations (NeurODEs)}~\cite{chen2018neural}, and an inviolable \textbf{3-second yellow clearance interval}, coupled with continuous-time discounting $\mathbf{\gamma^{\tau}}$.

\subsection{Summary of Contributions}
The primary contributions of this paper are summarized as follows:
\begin{itemize}
    \item \textbf{Formal SMDP Formulation with NeurODEs:} We formulate traffic signal actuation as a Semi-Markov Decision Process (SMDP) with continuous-time discounting $\mathbf{\gamma^{\tau}}$, continuous queue dynamics parameterized via \textbf{NeurODEs}~\cite{chen2018neural}, and explicit action-duration bounds ($\mathbf{10\,\text{s} \le \tau \le 45\,\text{s}}$, $\mathbf{t_{\text{yellow}} = 3\,\text{s}}$), eliminating dangerous phase flickering and ensuring compliance with municipal clearance standards.
    \item \textbf{Localized Rolling Emergency Preemption:} We develop a localized Rolling Emergency Preemption protocol that opens dynamic green waves along the emergency vehicle's immediate approach without paralyzing un-involved grid intersections.
    \item \textbf{Shockwave Mitigation and Fast Recovery ($T_{\text{rec}}$):} We introduce the Post-Preemption Shockwave Horizon ($T_{\text{rec}}$) and an active momentum-flushing reward mechanism, demonstrating that while contemporary preemption models induce acute cross-corridor gridlock, AURA actively absorbs and dissipates the secondary congestion shockwave \textbf{up to 5.6$\times$ faster}.
    \item \textbf{Multi-Baseline SOTA Empirical Benchmarking:} We conduct rigorous, multi-controller empirical evaluations within Eclipse SUMO under the Krauss car-following model against \textbf{five established baselines}: Traditional Fixed-Time (1958), Demand-Responsive Actuated control (1990), theoretical Max-Pressure (2013), modern SOTA PressLight (ACM KDD 2019), and contemporary EVP-DRL (Transportation Research Part C 2023).
    \item \textbf{Real-World Manhattan OpenStreetMap Scalability:} We validate the topological scalability of AURA on an authentic OpenStreetMap road network of \textbf{Midtown Manhattan} comprising \textbf{177 signalized traffic controllers} and \textbf{661 edges} under dense peak-hour flows (1,501 circulating vehicles), demonstrating genuine physical queue suppression and rapid emergency transit.
    \item \textbf{$\mathcal{O}(1)$ Decentralized Edge-Compute Efficiency:} We establish true $\mathcal{O}(1)$ decentralized complexity: requiring a parameter footprint of just \textbf{18.52\,KB} and single-core CPU inference latency of \textbf{33.28\,$\mu$s} (\textbf{30,050 decisions/sec/core}), demonstrating that a 10,000-intersection statewide network can be coordinated synchronously utilizing less than \textbf{33.3\% of a single commodity laptop CPU core} and under \textbf{20\,MB RAM} without server-farm infrastructure.
\end{itemize}

The remainder of this manuscript is organized as follows: Section~\ref{sec:related_work} reviews relevant literature in traffic control and reinforcement learning. Section~\ref{sec:methodology} presents the mathematical formulation and safety envelope mechanics. Section~\ref{sec:experiments} details the microscopic simulation setup, multi-baseline empirical findings, and recovery dynamics. Finally, Section~\ref{sec:conclusion} concludes the paper with directions for decentralized multi-agent deployment.
